\documentclass[9pt,a4paper]{extarticle}

% ATS-FRIENDLY AUTOMOTIVE RESUME
% Candidate: Vibhav Manikpuri | Target: Maruti Suzuki / Automotive R&D Internship

\usepackage[T1]{fontenc}
\usepackage[utf8]{inputenc}
\usepackage{lmodern}
\usepackage{microtype}
\usepackage[margin=0.42in]{geometry}
\usepackage{titlesec}
\usepackage{enumitem}
\usepackage{xcolor}
\usepackage[hidelinks]{hyperref}
\usepackage{ragged2e}

\definecolor{sectionblue}{RGB}{20,49,86}
\pagestyle{empty}
\raggedright
\setlength{\parindent}{0pt}
\setlength{\parskip}{0pt}
\setlist[itemize]{leftmargin=1.05em, itemsep=0.15pt, topsep=0.6pt, parsep=0pt, partopsep=0pt}
\renewcommand\labelitemi{\textbullet}

\titleformat{\section}{\vspace{-5pt}\normalsize\bfseries\color{sectionblue}}{}{0em}{}[\color{sectionblue}\titlerule\vspace{-6pt}]

\newcommand{\resumeSubheading}[4]{\vspace{1pt}\textbf{#1} \hfill \textbf{#2}\\[-2pt]\textit{#3} \hfill \textit{#4}\vspace{-2pt}}
\newcommand{\resumeProject}[3]{\vspace{1pt}\textbf{#1} \hfill \textbf{#2}\\[-2pt]\textit{#3}\vspace{-3pt}}
\newcommand{\skillcategory}[2]{\textbf{#1:} #2\\[-2pt]}

\begin{document}

\begin{center}
    {\Large \textbf{VIBHAV MANIKPURI}}\\[1pt]
    {\normalsize \textbf{Formula Student Vehicle Dynamics Engineer | Automotive Design Intern}}\\[1pt]
    Jabalpur, Madhya Pradesh, India \;|\; +91-9302752982 \;|\; \href{mailto:23bme070@iiitdmj.ac.in}{23bme070@iiitdmj.ac.in} \\
    \href{https://www.linkedin.com/in/vibhav-manikpuri-7a01192a2/}{LinkedIn} \;|\; \href{https://github.com/vibhavmanikpuri12-dev}{GitHub} \;|\; \href{https://github.com/vibhavmanikpuri12-dev/formula-student-vehicle-dynamics-portfolio}{Formula Student Portfolio}
\end{center}
\vspace{-9pt}

\section*{Professional Summary}
Final-year Mechanical Engineering undergraduate at PDPM IIITDM Jabalpur with hands-on Formula Student experience in vehicle dynamics, chassis subsystem design, CAD, CAE validation, manufacturing, and race-car integration. Worked on suspension, steering, hydraulic braking, upright, and wheel hub development using analytical calculations, SolidWorks, ANSYS Workbench, and engineering spreadsheets. Seeking an automotive internship in R\&D, product development, chassis design, vehicle testing, manufacturing engineering, or CAE.

\section*{Engineering Competencies}
\skillcategory{Vehicle Dynamics}{Suspension geometry, steering geometry, Ackermann calculation, brake system design, load transfer, roll center, ride rate, wheel rate, roll stiffness, vehicle packaging}
\skillcategory{CAD/CAE}{SolidWorks, CATIA V5, AutoCAD, ANSYS Workbench, static structural analysis, von Mises stress, deformation, factor of safety, fatigue, steady-state and transient thermal analysis}
\skillcategory{Manufacturing/Tools}{DFM, machining, welding coordination, grinding, sheet-metal work, assembly, fitment checks, troubleshooting, MS Excel, MATLAB, Python, C}

\section*{Automotive Engineering Experience}
\resumeSubheading{IIITDMJ Racing Formula Student Team - SAEINDIA SUPRA 2025}{2024 -- 2025}{Vehicle Dynamics Team Member, Cost Lead and Driver}{PDPM IIITDM Jabalpur}
\begin{itemize}
    \item Contributed to the end-to-end vehicle dynamics package of a Formula Student race car covering suspension, steering, hydraulic braking, upright, wheel hub, manufacturing, assembly, validation, and competition preparation.
    \item Converted analytical calculations and packaging constraints into manufacturable SolidWorks CAD models for steering, brake, upright, wheel hub, and suspension-related components.
    \item Supported fabrication and subsystem integration through machining coordination, grinding, welding support, fitment checks, assembly troubleshooting, and race-car preparation.
    \item Participated as driver and technical member, connecting vehicle-level feedback with steering, braking, and suspension design decisions during testing and competition preparation.
\end{itemize}

\section*{Selected Automotive Projects}
\resumeProject{Formula Student Suspension Geometry and Vehicle Dynamics Calculations}{Vehicle Dynamics}{Excel, analytical calculations, SolidWorks, manufacturing support}
\begin{itemize}
    \item Built a suspension calculation workflow for a 250 kg Formula Student vehicle using 1550 mm wheelbase, 1350/1250 mm front/rear track widths, 310 mm CG height, 45:55 weight distribution, and 1.5g lateral acceleration target.
    \item Calculated sprung/unsprung load distribution, vertical load transfer, rolling moment, ride rate, wheel rate, spring rate, sag, travel, roll center targets, ride frequency targets, and roll stiffness distribution.
\end{itemize}

\resumeProject{Formula Student Steering System Design}{Steering Kinematics}{Ackermann calculations, SolidWorks, manufacturing}
\begin{itemize}
    \item Designed steering geometry using Ackermann principles while balancing turning radius, tie-rod packaging, steering arm placement, chassis constraints, and manufacturability.
    \item Developed CAD models for steering components, coordinated manufacturing and assembly, and documented calculation steps, geometry references, component images, and validation notes.
\end{itemize}

\resumeProject{Formula Student Hydraulic Braking System}{Brake Design and CAE}{SolidWorks, ANSYS Workbench, thermal and structural FEA}
\begin{itemize}
    \item Engineered a hydraulic braking system for a 250 kg Formula Student vehicle by calculating static/dynamic weight transfer, required deceleration, braking force, brake torque, pedal ratio, master-cylinder pressure, and caliper clamping force.
    \item Sized and analyzed a 160 mm diameter, 3 mm thick SS410 brake rotor with 80 mm PCD; calculated 1:3 pedal ratio, 1500 N master-cylinder force, 5.26 MPa hydraulic pressure, and 3.81 kN clamping force.
    \item Performed ANSYS static structural, steady-state thermal, transient thermal, and brake pedal analyses to validate rotor safety, deformation, heat distribution, and pedal stiffness.
\end{itemize}

\resumeProject{Formula Student Upright and Wheel Hub Design}{Mechanical Design and CAE}{SolidWorks, ANSYS Workbench, fatigue analysis}
\begin{itemize}
    \item Designed front and rear uprights and wheel hubs considering suspension hard points, bearing fitment, brake mounting, wheel mounting, steering constraints, and manufacturability.
    \item Validated front hub, rear hub, front upright, and rear upright designs using ANSYS equivalent stress, total deformation, factor of safety, and fatigue analysis reports before manufacturing and integration.
\end{itemize}

\section*{Engineering Portfolio and GitHub}
\begin{itemize}
    \item Built a GitHub-ready Formula Student Vehicle Dynamics portfolio to organize subsystem calculations, CAD/CAE reports, manufacturing documentation, competition photographs, and resume evidence for recruiter review.
\end{itemize}

\section*{Education}
\resumeSubheading{PDPM Indian Institute of Information Technology, Design and Manufacturing, Jabalpur}{2023 -- 2027}{B.Tech in Mechanical Engineering}{Jabalpur, Madhya Pradesh}

\section*{Leadership and Achievements}
\begin{itemize}
    \item Vehicle Dynamics Cost Lead and active driver for SAEINDIA SUPRA 2025 Formula Student race-car development.
    \item Finalist, MP Government Hackathon - HACK \& MAKE 2025; contributed to engineering logic, cost modeling, scalability, and deployment planning.
    \item Additional projects: Exoskeleton Arm mechanical design in SolidWorks with servo actuation; AI-assisted Smart Irrigation System integrating sensors, actuators, and control logic.
\end{itemize}

\end{document}
